% GENERATED FILE. Edit canonical lessons, then run npm run generate:books.
% canonical-source-hash: fnv1a64:93dae4022a9774c3
% canonical-lessons: UR-C208-bulbul, UR-C208-saras, UR-C208-giddh, UR-C208-chamgadar, UR-C208-bichchhu

\chapter{Nightingale, Crane, Vulture, Bat, Scorpion}
\label{ch:ur-nightingale-crane-vulture-bat-scorpion}
\hlchaptermodality{208}

\begin{chapteropening}
\textbf{By the end of this chapter:} \emph{I can say a nightingale, a crane, a vulture, a bat and a scorpion.}

Complete the last lesson of chapter 208: I can say a nightingale, a crane, a vulture, a bat and a scorpion.
\end{chapteropening}

\section[\texorpdfstring{\ur{بلبل} (bulbul)}{بلبل (bulbul)}]{\ur{بلبل} (bulbul) — a nightingale}

\label{lesson:UR-C208-bulbul}



\begin{quote}
Before the new one: say the Urdu for black pepper, then the Urdu for vinegar.
\end{quote}

\subsection*{You'll want to know: \ur{بلبل}}
\textbf{\ur{بلبل}} — \emph{bulbul} — \textquotedblleft{}a nightingale\textquotedblright{}.

\begin{cousinweb}
\textbf{\ur{بلبل}} (\emph{bulbul}) is a loan, usually traced to Persian \emph{bulbul}, \textquotedblleft{}nightingale\textquotedblright{}.
\end{cousinweb}

\begin{grammarlens}[title={What you've built}]
One more word for this chapter.
\end{grammarlens}

\subsection*{Your turn}
\begin{itemize}
\raggedright
  \item \emph{Say it:} \emph{bulbul}
  \item \emph{Say it:} \emph{bulbul}, once more
  \item \emph{Say it:} say \emph{bulbul}
  \item \emph{Recall:} say the Urdu for black pepper, then the Urdu for vinegar, then say \emph{bulbul} again
\end{itemize}

\begin{tcolorbox}[breakable,colback=teal!4,colframe=teal!35!black,title={Before you move on}]
What does \ur{بلبل} mean? (\textquotedblleft{}A nightingale\textquotedblright{}.) Say it once more.
\end{tcolorbox}
\section[\texorpdfstring{\ur{سارس} (sāras)}{سارس (sāras)}]{\ur{سارس} (sāras) — a crane (bird)}

\label{lesson:UR-C208-saras}



\begin{quote}
Before the new one: say the Urdu for vinegar, then the Urdu for a nightingale.
\end{quote}

\subsection*{You'll want to know: \ur{سارس}}
\textbf{\ur{سارس}} — \emph{sāras} — \textquotedblleft{}a crane (bird)\textquotedblright{}.

\begin{cousinweb}
\textbf{\ur{سارس}} (\emph{sāras}) is a tadbhava, a word inherited from Sanskrit and reshaped by sound change. It comes from \emph{sārasa}, \textquotedblleft{}the crane\textquotedblright{} (CDIAL 13359). It is built on \emph{saras}, \textquotedblleft{}lake\textquotedblright{}.
\end{cousinweb}

\begin{grammarlens}[title={What you've built}]
One more word for this chapter.
\end{grammarlens}

\subsection*{Your turn}
\begin{itemize}
\raggedright
  \item \emph{Say it:} \emph{sāras}
  \item \emph{Say it:} \emph{sāras}, once more
  \item \emph{Say it:} say \emph{sāras}
  \item \emph{Recall:} say the Urdu for vinegar, then the Urdu for a nightingale, then say \emph{sāras} again
\end{itemize}

\begin{tcolorbox}[breakable,colback=teal!4,colframe=teal!35!black,title={Before you move on}]
What does \ur{سارس} mean? (\textquotedblleft{}A crane (bird)\textquotedblright{}.) Say it once more.
\end{tcolorbox}
\section[\texorpdfstring{\ur{گدھ} (giddh)}{گدھ (giddh)}]{\ur{گدھ} (giddh) — a vulture}

\label{lesson:UR-C208-giddh}



\begin{quote}
Before the new one: say the Urdu for a nightingale, then the Urdu for a crane.
\end{quote}

\subsection*{You'll want to know: \ur{گدھ}}
\textbf{\ur{گدھ}} — \emph{giddh} — \textquotedblleft{}a vulture\textquotedblright{}.

\begin{cousinweb}
\textbf{\ur{گدھ}} (\emph{giddh}) is a tadbhava, a word inherited from Sanskrit and reshaped by sound change. It comes from \emph{gṛdhra}, \textquotedblleft{}greedy; a vulture\textquotedblright{} (CDIAL 4233).
\end{cousinweb}

\begin{grammarlens}[title={What you've built}]
One more word for this chapter.
\end{grammarlens}

\subsection*{Your turn}
\begin{itemize}
\raggedright
  \item \emph{Say it:} \emph{giddh}
  \item \emph{Say it:} \emph{giddh}, once more
  \item \emph{Say it:} say \emph{giddh}
  \item \emph{Recall:} say the Urdu for a nightingale, then the Urdu for a crane, then say \emph{giddh} again
\end{itemize}

\begin{tcolorbox}[breakable,colback=teal!4,colframe=teal!35!black,title={Before you move on}]
What does \ur{گدھ} mean? (\textquotedblleft{}A vulture\textquotedblright{}.) Say it once more.
\end{tcolorbox}
\section[\texorpdfstring{\ur{چمگادڑ} (chamgādaṛ)}{چمگادڑ (chamgādaṛ)}]{\ur{چمگادڑ} (chamgādaṛ) — a bat (animal)}

\label{lesson:UR-C208-chamgadar}



\begin{quote}
Before the new one: say the Urdu for a crane, then the Urdu for a vulture.
\end{quote}

\subsection*{You'll want to know: \ur{چمگادڑ}}
\textbf{\ur{چمگادڑ}} — \emph{chamgādaṛ} — \textquotedblleft{}a bat (animal)\textquotedblright{}.

\begin{cousinweb}
\textbf{\ur{چمگادڑ}} (\emph{chamgādaṛ}) goes back to Sanskrit \emph{carmacaṭaka}, \textquotedblleft{}bat\textquotedblright{}, \textquotedblleft{}the leather sparrow\textquotedblright{}, built from \emph{carman}, \textquotedblleft{}skin, leather\textquotedblright{}, and \emph{caṭaka}, \textquotedblleft{}sparrow\textquotedblright{}. CDIAL lists Hindi \emph{cām-gīdaṛ}, \textquotedblleft{}bat\textquotedblright{}, under it (CDIAL 4700).
\end{cousinweb}

\begin{grammarlens}[title={What you've built}]
One more word for this chapter.
\end{grammarlens}

\subsection*{Your turn}
\begin{itemize}
\raggedright
  \item \emph{Say it:} \emph{chamgādaṛ}
  \item \emph{Say it:} \emph{chamgādaṛ}, once more
  \item \emph{Say it:} say \emph{chamgādaṛ}
  \item \emph{Recall:} say the Urdu for a crane, then the Urdu for a vulture, then say \emph{chamgādaṛ} again
\end{itemize}

\begin{tcolorbox}[breakable,colback=teal!4,colframe=teal!35!black,title={Before you move on}]
What does \ur{چمگادڑ} mean? (\textquotedblleft{}A bat (animal)\textquotedblright{}.) Say it once more.
\end{tcolorbox}
\section[\texorpdfstring{\ur{بچھو} (bichchhū)}{بچھو (bichchhū)}]{\ur{بچھو} (bichchhū) — a scorpion}

\label{lesson:UR-C208-bichchhu}



\begin{quote}
Before the new one: say the Urdu for a vulture, then the Urdu for a bat.
\end{quote}

\subsection*{You'll want to know: \ur{بچھو}}
\textbf{\ur{بچھو}} — \emph{bichchhū} — \textquotedblleft{}a scorpion\textquotedblright{}.

\begin{cousinweb}
\textbf{\ur{بچھو}} (\emph{bichchhū}) is a tadbhava, a word inherited from Sanskrit and reshaped by sound change. It comes from \emph{vṛścika}, \textquotedblleft{}scorpion\textquotedblright{} (CDIAL 12081). CDIAL asks whether the many shapes of this word come from a taboo on naming the scorpion.
\end{cousinweb}

\begin{grammarlens}[title={What you've built}]
One more word for this chapter.
\end{grammarlens}

\subsection*{Your turn}
\begin{itemize}
\raggedright
  \item \emph{Say it:} \emph{bichchhū}
  \item \emph{Say it:} \emph{bichchhū}, once more
  \item \emph{Say it:} say \emph{bichchhū}
  \item \emph{Recall:} say the Urdu for a vulture, then the Urdu for a bat, then say \emph{bichchhū} again
\end{itemize}

\begin{tcolorbox}[breakable,colback=teal!4,colframe=teal!35!black,title={Before you move on}]
What does \ur{بچھو} mean? (\textquotedblleft{}A scorpion\textquotedblright{}.) Say it once more.
\end{tcolorbox}
